%=== CHAPTER ONE (1) ===
%=== INTRODUCTION ===

\chapter{Introduction}
\label{ch:introduction}

\begin{spacing}{2.0}
%\setlength{\parskip}{0.2in}

\section{Background and Context}
The rapid and sustained expansion of global urban infrastructure has resulted in an unprecedented increase in vehicular traffic. Consequently, this heightened density has led to a greater frequency of road traffic collisions and systemic congestion, thereby presenting significant risks to public safety and economic efficiency. In response to these mounting challenges, contemporary municipal authorities are increasingly adopting Intelligent Transport Systems (ITS) to safeguard public infrastructure and optimise traffic management. Historically, traffic monitoring has relied heavily on intrusive hardware, such as inductive loop detectors, or manual human observation. However, these traditional methods exhibit notable operational limitations; physical sensors are unable to provide comprehensive spatial context, while human observation is inherently susceptible to fatigue-induced errors and delayed response times. To overcome these deficiencies, the incorporation of deep learning and computer vision technologies into existing Closed-Circuit Television (CCTV) networks has emerged as a highly effective, non-intrusive approach for collecting high-dimensional data related to vehicle trajectories, interactions, and anomalies.

\section{Problem Statement}
Notwithstanding these advancements, the task of robust traffic anomaly detection remains a highly complex challenge within the domain of computer vision. Contemporary object detection frameworks, although markedly effective at identifying individual vehicles under optimal conditions, frequently prove insufficient when confronted with the chaotic and unstructured spatial dynamics inherent in physical collisions. During such incidents, the visual representations of vehicles tend to intersect, deform, or fragment, leading to the failure of traditional bounding-box-based spatial tracking methods to sustain consistent identification.

Conversely, although sophisticated spatiotemporal networks, such as Transformer-based architectures, possess the capability to model these interactions effectively over time, they demand substantial computational resources. This considerable processing requirement renders them largely impractical for real-time deployment on resource-constrained edge devices at municipal level. Systems reliant on inflexible spatial tracking or highly complex temporal models often exhibit unacceptable false-positive rates or experience debilitating processing delays, thereby compromising their suitability for life-saving early warning applications.

\section{The Proposed Hybrid Architecture}
To address this technological disparity, the present research introduces an innovative, solely spatiotemporal, vision-based framework specifically designed for real-time traffic anomaly detection. The primary advancement of this architecture lies in the complete abandonment of fragile spatial vehicle tracking, in favour of assessing unstructured kinetic energy.

The system utilises an asynchronous, dual-threaded producer-consumer architecture designed to optimise hardware concurrency while avoiding stalling during video ingestion. The primary thread continuously evaluates the physical motion within the scene using the Farnebäck algorithm for dense optical flow. To minimise environmental noise, stringent programme thresholds and Region of Interest (ROI) masks are implemented. This raw kinetic data is subsequently encoded into a spatiotemporal Hue-Saturation-Value (HSV) tensor. Once generated, these tensors are classified asynchronously via a customised ResNet-18 Convolutional Neural Network (CNN), serving as the probabilistic verification mechanism to accurately differentiate between routine traffic flow and catastrophic structural impacts.


\section{Aims and Objectives}
The primary objective of this research is to develop a highly precise and computationally efficient prototype capable of maintaining a throughput of sixty frames per second on standard consumer-grade hardware. The system seeks to utilise the deterministic physics of dense optical flow in conjunction with the probabilistic pattern recognition provided by deep learning.

In pursuit of this aim, the network was systematically trained and evaluated utilizing a mathematically balanced synthesis of the Car Accident Detection and Prediction (CADP) and Traffic Anomaly Detection (TAD) benchmark datasets. To address the inherent class imbalance characteristic of traffic data, a Binary Cross-Entropy (BCE) loss function was adopted, thereby compelling the network to dynamically refine the boundary between complex, non-anomalous kinetic events and genuine collisions. Ultimately, the objective is to attain a near-instantaneous detection latency for Time-to-Accident (TTA), thus establishing a highly responsive and privacy-conscious foundation for automated municipal alert systems.

\section{Research Questions}
To evaluate the efficacy and practical applicability of the proposed purely spatiotemporal architecture, this study systematically addresses the following primary research questions:
\begin{itemize}
    \item \textbf{RQ1:} Is it possible to effectively detect traffic abnormalities, such as multi-vehicle collisions, using a purely spatiotemporal vision-based architecture that relies exclusively on dense optical flow and deep residual networks, thereby bypassing traditional object detection?
    \item \textbf{RQ2:} What are the primary computational and environmental challenges associated with evaluating pure kinetic energy for anomaly detection, and how does the proposed architecture mitigate these vulnerabilities?
\end{itemize}

\section{Structure of the Thesis}
The remainder of this thesis is organised as follows. Chapter 2 provides a comprehensive Literature Review, examining the development of object detection, multi-object tracking, and spatiotemporal optical flow analysis within the context of Intelligent Transportation Systems (ITS). Chapter 3 details the Methodology, outlining the software architecture, data engineering pipelines, mathematical transformations, and the dual-threaded integration of the spatiotemporal system. Chapter 4 presents the Analysis, Results, and Discussion, wherein the framework is evaluated using quantitative classification metrics as well as qualitative real-world constraints. Finally, Chapter 5 concludes the study by summarising the principal contributions to knowledge, recognising the inherent limitations of the system, and offering strategic recommendations for future development.

\end{spacing}
%=== END OF CHAPTER ONE ===
\newpage


